\section{Background}
\label{sec:background}

\subsection{Total variation and what a bound buys}
\label{sec:background/tv}

For two distributions $P$ and $Q$ over the same finite alphabet, the total
variation distance $\tv(P,Q)$ in \eqref{eq:tv} is, up to a factor of one half,
the smallest error probability achievable by any test that must decide which of
the two generated an observation. This is why a bound on $\tv$ is the natural
quantity for a practitioner. An upper bound on $\tv$ translates into a floor
that no classifier can go below; a lower bound translates into a statement that
the pair is genuinely separable. Neither is available from a classifier's own
accuracy, which measures one system rather than the problem.

We write the bracket as in \eqref{eq:bracket} and call each endpoint by what it
is. The upper endpoint is finite-sample and holds at the observed flow counts
with a stated confidence, under no assumption on $P$ or $Q$ beyond a finite
alphabet; the lower endpoint is a classifier-based bound at a nominal level,
distribution-free in the same sense but conditional on the fitted decision
function. Both carry a confidence level, so we report them separately and quote
the gate as a conjunction \eqref{eq:gate} rather than folding them into one
number.

\subsection{Prior information bounds on traffic analysis}
\label{sec:background/prior}

Bounding traffic-analysis performance from information-theoretic quantities,
rather than measuring it, is well established. \citet{cherubin2017bayes} gives
feature-set-conditional Bayes-error bounds for website fingerprinting, a
closed-set floor and the closest precedent for our upper side. Work on anonymity
and quantitative information flow follows the same route, from an entropy-like
quantity or a metric on distributions toward a bound on inference
\citep{serjantov2003metric,ghaderi2010anonymous,venkitasubramaniam2008eavesdroppers,venkitasubramaniam2008chaum,chatzikokolakis2007error,chatzikokolakis2008noisy}.
Our construction differs in being an explicit interval on real trace data, with
the difference between its endpoints kept visible rather than averaged away.

The most directly relevant recent results are the side-channel leakage
inevitability theorem of \citet{liu2026inevitability} and its rate-distortion
companion \citep{liu2026rate}. The former establishes that leakage is
unavoidable under explicit non-degeneracy and distinguishability conditions; the
latter characterizes an exact leakage-cost frontier. Both are single-letter in a
specific sense. The exponential error decay in the first requires conditionally
i.i.d.\ repeated observations, and the frontier in the second is proved within a
stationary-memoryless class. Packet prefixes satisfy neither condition. Our
measurement makes no such assumption and instead reports what is certifiable
about the dependent prefix law as it stands, which is why the upper side is
distribution-free, while the lower side is a classifier-based bound whose
width is set by the adversarial choice of decision function rather than by an
exponential decay rate.

\subsection{Estimating and testing total variation}
\label{sec:background/tv_est}

Turning a bound like \eqref{eq:upper} into a usable statement requires a
confidence radius for the empirical distribution, and the classical
Kolmogorov-Smirnov style constructions supply exactly that
\citep{wei2011dkwtwo,reeve2024dkwmassart}, with high-probability and two-sample
refinements that make the constant tight enough to be informative
\citep{diakonikolas2020highprob,han2015l1distribution,jiao2018l1distance}. Two
further routes are relevant. Classifier-based methods give lower bounds from a
decision function's error rate, and \citet{michel2020high} provides the adaptive
high-probability form we use for the lower side. \citet{tao2024discriminative}
relates discriminative estimation of total variation to classification risk
under a hypothesis-class assumption, which is the route we take for the
conditional large-prefix bound where the plug-in certificate is unavailable.

Open-set generalization and learnability results
\citep{fang2021learning,garov2025closer} bound the risk of a learning procedure;
ours bounds the separation available in the data, and the two are complementary
since a learnability bound says how well a detector can be trained while a
bracket says how much room there is to train into. The unknown-class mass a
bracket over observed services omits is a separate problem with its own
literature \citep{orlitsky2016optimal,eriksson2026unseen}, used for the
site-level term in \Cref{sec:discussion/missing_mass}.

\subsection{The QUIC observable}
\label{sec:background/quic}

QUIC's wire image exposes a packet-level view from which an observer recovers
packet directions, sizes, and timing, while content stays behind encryption;
\citet{kuhlewind2022rfc} documents this for operators. We use the
(direction, size) sequence of the leading packets, with prefix length and
alphabet resolution as the two parameters we pre-register and vary. Timing is
deliberately excluded, and we measure the cost of excluding the host identifier
rather than asserting it is free.
